\documentclass[10pt,a4paper]{amsart}
\usepackage[margin=19mm]{geometry}
\usepackage{amsmath,amssymb,mathtools}
\usepackage[scr=rsfs,cal=euler]{mathalfa}
\usepackage{xfrac}
\usepackage[unicode,hidelinks,bookmarks=false]{hyperref}
\newcommand{\ie}{i.e.\ }
\newcommand{\cf}{cf.\ }
\newcommand{\resp}{respectively\ }
\newcommand{\Subsection}{\subsection}
\providecommand{\nobreakdash}{\nobreak\hbox{-}}
\setlength{\emergencystretch}{2em}
\multlinegap0pt
\usepackage{tikz}
\usetikzlibrary{cd}
\usepackage{amssymb}
\usepackage[shortlabels]{enumitem}
\usepackage{xcolor}
\usepackage{mathtools}
\newtheorem{proposition}{Proposition}
\newcommand{\E}{\mathcal{E}}
\DeclareMathOperator{\Hom}{Hom}
\DeclareMathOperator{\Inn}{Inn}
\newcommand{\PAlg}{\ensuremath{\mathsf{PAlg}}}
\DeclareMathOperator{\SplExt}{SplExt}
\newcommand{\V}{\mathcal{V}}
\title{On the representability of actions\\of unital algebras}
\author{Manuel Mancini, Federica Piazza, Corentin Vienne}
\date{Source: arXiv:2503.04488v2 (CC BY 4.0)}
\begin{document}
\maketitle

\noindent\emph{Source and licence.} On the representability of actions\\of unital algebras. Version arXiv:2503.04488v2 (CC BY 4.0), distributed by arXiv under the Creative Commons Attribution 4.0 International licence, http://creativecommons.org/licenses/by/4.0/, as recorded in the arXiv licence metadata for this exact version and shown on the abstract page https://arxiv.org/abs/2503.04488v2. Full licence notice: \texttt{licenses/arxiv-2503.04488-CC-BY-4.0.txt}.\par\smallskip

\noindent\textit{Editorial scope.} Counted unit: the article's proposition prop\_e(x) (source0.tex lines 530--532). The article's complete original proof of it is delivered with it (source0.tex lines 534--570, one \texttt{proof} environment of 37 lines). No mathematics is written, repaired or re-derived: the statement and the proof are the article's own lines, in the article's own notation, and no retained line is substituted or re-flowed.  No further source-local context is needed: every cross-reference of every retained line resolves inside the two delivered spans.  Nothing was omitted from any retained span, so the package carries no omission record.  No retained line contains a citation, so the wrapper carries no bibliography and no bibliography entry.  The retained lines contain the article's own diagram code, which the wrapper typesets with the standard tikz packages; no image file is delivered and the package declares no asset dependency.  Editorial formatting only, in addition: an AMS \texttt{amsart} preamble that restates the article's own declarations and macros used by the retained lines, namely the environments \texttt{proposition} on separate counters, together with the standard packages those lines need.  The retained environments print in this document as Proposition 1 in order of appearance, whereas the article prints the same items with the numbering of its own full document.  The complete native source of the article as distributed in the arXiv e-print for this exact version is bundled unchanged as \texttt{sources/174333/source0.tex}.
\begin{proposition}\label{prop_e(x)}
Let $\V$ be an operadic, action accessible variety of non-associative algebras, and fix a choice of $\lambda/\mu$-rules. For any object $X \in \V$, if the partial algebra homomorphism $\Inn \colon \widetilde{U}(X) \to \E(X)$ is an isomorphism, then the actions on $X$ are representable and $X$ is its own actor.
\end{proposition}

\begin{proof}
Suppose that $\Inn$ is an isomorphism. Then
\[
\Hom_{\PAlg}(\widetilde{U}(-),\E(X))=\Hom_{\V}(-,\E(X)) \cong \Hom_{\V}(-,X)
\]
and hence, by composition, we get a natural monomorphism of functors
\[
\widetilde{\tau} \colon \SplExt(-,X) \rightarrowtail \Hom_{\V}(-,X).
\]
It remains to show that, for any algebra $B$ of $\V$, the component $\widetilde{\tau}_B$ is surjective, that is, every homomorphism $\varphi \colon B \to X$ in $\V$ arises from a split extension of $B$ by~$X$ in $\V$.

To this end, consider the split extension
\begin{equation}\label{split_phi}
\begin{tikzcd}
X \arrow[r, "\iota_2"] &
B \ltimes X \arrow[r, shift left, "\pi_1"] &
B, \arrow[l, shift left, "\iota_1"]
\end{tikzcd}
\end{equation}
where $B \ltimes X$ is the algebra whose underlying vector space is $B \times X$ and whose multiplication is defined by
\[
(b,x) \cdot (b',x') = (bb', xx' + \varphi(b)x' + x\varphi(b')),
\]
and $\iota_1, \iota_2$ and $\pi_1$ are the canonical injections and projection. Straightforward computations show that $\tau_B$ sends \eqref{split_phi} to the morphism
\[
\Inn \circ \varphi \colon B \to \E(X) \colon b \mapsto (L_{\varphi(b)}, R_{\varphi(b)}).
\]
Consequently, $\widetilde{\tau}_B$ sends the split extension \eqref{split_phi} to the morphism $\varphi$.

Finally, we need to check that \eqref{split_phi} is indeed a split extension in $\V$. Let $\Phi(x_1,\ldots,x_k)=0$ be any multilinear identity of $\V$. Then, for any $\alpha_1,\ldots,\alpha_k \in B \cup X$, the evaluation $\Phi(\alpha_1,\ldots,\alpha_k)$ reduces to $0$, since
\[
b \ast x = \varphi(b)x \; \text{ and } \; x \ast b = x\varphi(b),
\]
for every $b \in B$ and $x \in X$. Thus, any such evaluation reduces to an expression involving only multiplications in $X$, and hence vanishes. It follows that $B \ltimes X$ satisfies all identities of $\V$, and therefore is an algebra of $\V$.

Therefore, $\widetilde{\tau}_B$ is surjective, and the actions on $X$ are representable. Since $\Inn$ is an isomorphism, it follows that $X$ is its own actor.
\end{proof}

\end{document}
